

%While the first part of my work addresses inference given limited data, the second focuses on how to strategically acquire new data to overcome these limits. In this line of work, I aimed to develop methods to collect most informative or alternative data---taking into account not only the data themselves but also prior knowledge that helps the inference, while considering robustness---we need to know if a small amount of data determined conclusion.  

%In a complementary line of work, I develop strategies for collecting new and informative data when existing observations are limited or noisy. These methods incorporate prior knowledge to guide adaptive experimental design, aiming to identify which data would most improve inference and to assess how robust conclusions remain under limited sampling. Together, they close the loop between inference and data acquisition.

%In a complementary line of work, I develop strategies for data acquisition and reasoning under limited data---both by collecting the most informative new data and by assessing the robustness of inferences drawn from existing data. These methods close the loop between inference, data acquisition, and reliability assessment, asking not only what data we should collect next but also how confident we can be in the conclusions supported by the data we already have.

Beyond the midstream of the inference pipeline, I study both upstream and downstream processes—data acquisition and robustness assessment—as two sides of the same problem: determining what information matters most. When data are limited, we face two complementary challenges: prospectively, deciding which new, unseen measurements will be most informative; and retrospectively, assessing how much our conclusions depend on the limited data already observed. I develop methods that close the loop between inference, data design, and reliability evaluation, ensuring that what we learn is both informative and stable.